\documentclass[10pt]{article}
\usepackage[T1]{fontenc}
\usepackage[utf8]{inputenc}
\usepackage[french]{babel}
\usepackage{amsmath,amssymb,mathrsfs,mathpazo}
\usepackage{geometry}
\geometry{a4paper,margin=2cm}
\usepackage{graphicx,array,tabularx,tabularray,colortbl,multirow,multicol,booktabs,enumitem,pifont,ulem,fancybox,diagbox,cancel}
\usepackage{pstricks,pstricks-add,pst-plot,pst-tree,pst-node,pst-text,pst-func,pst-math,pst-bspline,pst-3dplot,pst-eucl,pst-solides3d}
\let\pscalculate\undefined
\usepackage{pst-bezier}
\usepackage{tikz,pgfplots,tkz-tab}
\usetikzlibrary{arrows,arrows.meta,patterns,calc,positioning,trees,decorations.markings,shapes.geometric,shapes.misc,shapes.arrows,intersections,decorations.pathmorphing,decorations.pathreplacing,backgrounds,shadows,babel}
\pgfplotsset{compat=1.18}
\usepgfplotslibrary{fillbetween}
\usepackage[np]{numprint}
\usepackage{listings,makecell,pas-tableur,twoopt}
\let\barre\undefined
\usepackage[pstricks]{tablvar}
\definecolor{darkgreen}{RGB}{0,100,0}
\providecommand{\red}{\color{red}}
\providecommand{\blue}{\color{blue}}
\providecommand{\green}{\color{green}}
\providecommand{\black}{\color{black}}
\providecommand{\white}{\color{white}}
\providecommand{\gray}{\color{gray}}
\providecommand{\cyan}{\color{cyan}}
\providecommand{\magenta}{\color{magenta}}
\providecommand{\yellow}{\color{yellow}}
\providecommand{\violet}{\color{violet}}
\providecommand{\orange}{\color{orange}}
\providecommand{\euro}{\texteuro}
\providecommand{\EUR}{\texteuro}
\providecommand{\ieme}{\textsuperscript{e}}
\providecommand{\degres}{\ensuremath{^\circ}}
\providecommand{\degree}{\ensuremath{^\circ}}
\providecommand{\vv}[1]{\overrightarrow{#1}}
\providecommand{\Coord}[1]{#1}
\providecommand{\decosix}{}
\providecommand{\starredbullet}{\ensuremath{\bullet}}
\pagestyle{empty}
\setlength{\parindent}{0pt}
\usepackage{adjustbox}\DeclareUnicodeCharacter{2236}{:}\DeclareUnicodeCharacter{2190}{\ensuremath{\leftarrow}}
\providecommand{\R}{\mathbb{R}}
\providecommand{\N}{\mathbb{N}}
\providecommand{\Z}{\mathbb{Z}}
\providecommand{\Q}{\mathbb{Q}}
\providecommand{\D}{\mathbb{D}}
\providecommand{\C}{\mathbb{C}}
\providecommand{\e}{{\mathrm{e}}}
\providecommand{\vect}[1]{\overrightarrow{#1}}
\providecommand{\vectt}[1]{\overrightarrow{\mathrm{#1}}}
\providecommand{\vvt}[1]{\overrightarrow{\mathrm{#1}}}
\providecommand{\np}[2][]{{#2~#1}}
\providecommand{\vv}[1]{{\overrightarrow{#1}}}
\renewcommand{\d}{{\mathrm{d}}}
\providecommand{\degres}{{^\circ}}
\providecommand{\degre}{{^\circ}}
\providecommand{\blue}{\color{blue}}
\providecommand{\red}{\color{red}}
\providecommand{\enskip}{\quad}
\providecommand{\llbracket}{\lbrack\!\lbrack}
\providecommand{\rrbracket}{\rbrack\!\rbrack}
\providecommand{\euro}{{\text{€}}}
\providecommand{\no}{{\text{n°}}}
\providecommand{\interval}[2]{{[#1;#2]}}
\providecommand{\intervFO}[2]{{[#1;#2[}}
\providecommand{\dotfill}{\ldots}

\providecommand{\ds}{\displaystyle}
\providecommand{\textding}[1]{#1}
\providecommand{\Oij}{\(\mathrm{O};\vec{\imath},\vec{\jmath}\)}
\providecommand{\Oijk}{\(\mathrm{O};\vec{\imath},\vec{\jmath},\vec{k}\)}
\providecommand{\Ouv}{\(\mathrm{O};\vec{u},\vec{v}\)}
\providecommand{\intervFF}[2]{\ensuremath{[#1;#2]}}
\renewcommand{\euro}{\ifmmode\text{\texteuro}\else\texteuro\fi}
\DeclareUnicodeCharacter{00B0}{\ifmmode^{\circ}\else\textdegree\fi}
\DeclareUnicodeCharacter{20AC}{\ifmmode\text{\texteuro}\else\texteuro\fi}
% Énoncés : Éducation nationale. Transcription et corrigés : APMEP, Denis Vergès et contributeurs. Découpage et classement : CoopMaths, Sébastien Lozano, Éric Elter et contributeurs. Adaptation de la présentation : Maths Explorer.
% CC BY-SA : https://coopmaths.fr/www/about/
% https://coopmaths.fr/alea/static/eam/2026/tex/eam_specifique_2026_06_polynesie_1.tex
\begin{document}
\noindent\textit{APMEP / CoopMaths. CC BY-SA. Présentation adaptée par Maths Explorer.}\par\medskip
Deux entreprises A et B souhaitent réduire leurs déchets plastiques.

\begin{itemize}[label={\textbullet~}]
\item L'entreprise A réduit ses déchets de 15 tonnes chaque année.

Elle produit 500 tonnes de déchets en 2020.

\item L'entreprise B réduit ses déchets de $8\,\%$ chaque année.

Elle produit 600 tonnes de déchets en 2020.
\end{itemize}

\subsection*{Partie A}
On a représenté ci-dessous les quantités annuelles estimées de déchets plastiques des entreprises A et B entre 2020 et 2030, en prenant l'année 2020 comme année de rang 0 .

\begin{center}
\psset{xunit=1cm,yunit=0.01cm,arrowsize=2pt 3}
\begin{pspicture*}(-1,-20)(10.5,750)
\psaxes[linewidth=1.25pt,labelFontSize=\scriptstyle,Dy=100]{->}(0,0)(0,0)(10.5,750)
\multido{\n=0+1}{11}{\psline[linewidth=0.15pt](\n,0)(\n,700)}
\multido{\n=0+100}{8}{\psline[linewidth=0.15pt](0,\n)(10,\n)}
\psdots(0,600)(1,550)(2,505)(3,475)(4,430)(5,398)(6,370)(7,340)(8,300)(9,285)(10,270)
\psdots[dotstyle=+,dotangle=45,dotscale=1.5,linewidth=1.5pt](0,500)(1,485)(2,470)(3,455)(4,440)(5,425)(6,410)(7,395)(8,380)(9,365)(10,350)
\uput[u](9,0){Rang de l'année}\uput[r](0,730){Quantité de déchets (en tonnes)}
\end{pspicture*}
\end{center}

\medskip

\begin{enumerate}
  \item Justifier que la quantité de déchets produits par l'entreprise A est représentée par les points marqués d'une croix.
  \item Déterminer graphiquement à partir de quelle année l'entreprise $B$ produit moins de déchets que l'entreprise A.
  \item L'entreprise A atteindra-t-elle son objectif de réduire au moins de moitié ses déchets plastiques en 2030 ?
\end{enumerate}

\subsection*{Partie B}

\begin{enumerate}
\item On note $u_{n}$ la quantité de déchets plastiques, en tonnes, produits par l'entreprise A lors de l'année $2020+n$, où $n$ est un entier naturel.

Ainsi,\quad $u_{0}=500$.
	\begin{enumerate}
		\item Calculer $u_{1}$.

		Interpréter ce résultat dans le contexte de l'exercice.

		\item Justifier que la suite $(u_{n})$ est arithmétique.

		Préciser sa raison.

	\end{enumerate}
		On a utilisé un tableur pour calculer les termes de la suite $(u_{n})$.

		Un extrait de la feuille de calcul est donné ci-dessous.

\begin{center}

\begin{tabular}{|l|c|c|c|}
\hline
 & A & B & C \\
\hline
1 & \multicolumn{3}{|c|}{Déchets produits par l'entreprise A} \\
\hline
2 & Année & Rang $n$ & Quantité de déchets $u_{s}$ (en tonnes) \\
\hline
3 & 2020 & 0 & 500 \\
\hline
4 & 2021 & 1 & 485 \\
\hline
5 & 2022 & 2 & 470 \\
\hline
6 & 2023 & 3 & 455 \\
\hline
7 & 2024 & 4 & 440 \\
\hline
8 & 2025 & 5 & 425 \\
\hline
9 & 2026 & 6 & 410 \\
\hline
10 & 2027 & 7 & 395 \\
\hline
11 & 2028 & 8 & 380 \\
\hline
12 & 2029 & 9 & 365 \\
\hline
13 & 2030 & 10 & 350 \\
\hline
14 & 2031 & 11 & 335 \\
\hline
15 & 2032 & 12 & 320 \\
\hline
16 & 2033 & 13 & 305 \\
\hline
17 & 2034 & 14 & 290 \\
\hline
18 & 2035 & 15 & 275 \\
\hline
19 & 2036 & 16 & 260 \\
\hline
20 & 2037 & 17 & 245 \\
\hline
21 & 2038 & 18 & 230 \\
\hline
22 & 2039 & 19 & 215 \\
\hline
23 & 2040 & 20 & 200 \\
\hline
\end{tabular}
\end{center}
	\begin{enumerate}[resume]
		\item  Dans la cellule \textsf{C4}, quelle formule doit-on saisir puis recopier vers le bas, afin d'obtenir les termes de la suite $\left(u_{n}\right)$ ?
		\item  Déterminer au bout de combien d'années l'entreprise A atteindra son objectif de réduire au moins de moitié ses déchets plastiques.
	\end{enumerate}
\item On note $v_{n}$ la quantité de déchets plastiques, en tonnes, produits par l'entreprise B lors de l'année $2020+n$, où $n$ est un entier naturel.

Ainsi, \quad $v_{0}=600$.
	\begin{enumerate}
		\item Justifier que pour tout entier naturel $n$,
		\[v_{n+1}=0,92 v_{n}.\]

		\item La suite $\left(v_{n}\right)$ est-elle arithmétique ou géométrique?

		Préciser sa raison.

		\item Exprimer $v_{n}$ en fonction de $n$.

		\item Au bout de combien d'années l'entreprise B aura-t-elle atteint son objectif de réduire au moins de moitié ses déchets plastiques ?

		Justifier votre réponse à l'aide du tableau ci-dessous.

			\begin{tblr}{|X[l,2.5cm,m]|*{10}{X[c]|}}
				\hline
				$n$ & 1 & 2 & 3 & 4 & 5 & 6 & 7 & 8 & 9 & 10 \\
				\hline
				\small $0,92^{n}$ (arrondi au
				centième)& 0,92 & 0,85 & 0,78 & 0,72 & 0,66 & 0,61 & 0,56 & 0,51 & 0,47 & 0,43 \\
				\hline
			\end{tblr}

	\end{enumerate}
\end{enumerate}

\end{document}
